
\documentclass[11pt,a4paper]{article}

\usepackage[utf8]{inputenc}
\usepackage[T1]{fontenc}
\usepackage[spanish]{babel}

\usepackage{tgpagella}
\usepackage{geometry}
\usepackage{graphicx}
\usepackage{xcolor}
\usepackage{titlesec}
\usepackage{fancyhdr}
\usepackage[parfill]{parskip}
\usepackage[expansion=false]{microtype}
\usepackage{hyperref}
\usepackage{setspace}
\usepackage{needspace}

\geometry{top=3.0cm,bottom=3.0cm,left=3.5cm,right=3.5cm}

\setlength{\parskip}{0.65em}
\setlength{\parindent}{0em}

\definecolor{azulai}{RGB}{30,80,140}
\definecolor{doradoai}{RGB}{140,90,0}
\definecolor{grisai}{RGB}{70,70,70}

\titleformat{\section}
{\Large\bfseries\color{azulai}}
{}{0em}{}
[{\color{azulai}\titlerule[0.5pt]}]

\titlespacing*{\section}{0pt}{1.8em}{0.8em}

\titleformat{\subsection}
{\large\bfseries\color{doradoai}}
{}{0em}{}

\titlespacing*{\subsection}{0pt}{1.2em}{0.4em}

\pagestyle{fancy}
\fancyhf{}

\rhead{\textcolor{grisai}{\small Izas — Arquitectura Interna}}
\lhead{\textcolor{grisai}{\small Carta Natal Base}}

\cfoot{\textcolor{grisai}{\small\thepage}}

\renewcommand{\headrulewidth}{0.3pt}

\hypersetup{
colorlinks=true,
linkcolor=azulai,
urlcolor=azulai
}

\setstretch{1.4}

\begin{document}

% ── PORTADA ──────────────────────────────────────────────────────────────

\begin{titlepage}

\centering

\vspace*{1.5cm}

{\Huge\bfseries\color{azulai} Carta Natal Base}\\[0.5cm]

{\large\color{grisai} Arquitectura Interna}\\[0.4cm]

{\small\itshape\color{grisai}
Un método para sostener cuerpo, energía y vida con coherencia
}\\[2cm]

{\huge\color{doradoai} Izas}\\[1.5cm]

{\Large 16/04/1979 \quad 04:20}\\[0.3cm]

{\Large Pamplona, España}\\[0.3cm]

{\normalsize
Lat: 42.8157° \quad
Lon: -1.6522° \quad
UTC+02:00
}\\[0.5cm]

{\normalsize
Ascendente: Acuario 6°20'
}\\[2.5cm]

\vfill

{\small Generado el 21/05/2026}

\end{titlepage}


% ── INTRODUCCIÓN ─────────────────────────────────────────────────────────────

\section*{Introducción}

Esta carta natal base no busca definir quién eres.

La astrología se utiliza aquí como lenguaje de observación:
una forma de mirar cómo se organiza tu energía,
qué registros aparecen con más facilidad
y qué dinámicas pueden influir en tu manera de vivir, sostenerte y relacionarte con el entorno.

No se trata de una lectura predictiva ni de una identidad fija.
Es un mapa inicial de orientación.

Las siguientes páginas muestran únicamente las capas principales de la carta:
los ejes más visibles,
la distribución general de la energía
y algunos patrones básicos desde los que sueles moverte.

\vspace{0.4cm}

\newpage


% ── RUEDA ────────────────────────────────────────────────────────────────

\section{Carta Natal}

\begin{center}
\includegraphics[width=0.8\textwidth]{Izas_carta_base_rueda.png}
\end{center}

\vspace{0.8cm}

% ── POSICIONES PRINCIPALES ───────────────────────────────────────────────

\section{Posiciones principales}

\begin{center}

Sol en Aries — Casa 2\\
Luna en Sagitario — Casa 10\\
Ascendente Acuario\\
Medio Cielo en Escorpio

\end{center}

\vspace{0.5cm}

{\small\itshape
La lectura ampliada desarrolla el mapa completo de posiciones,
aspectos y configuraciones de la carta.
}

\vspace{0.8cm}
\Needspace{5\baselineskip}

% ── VISIÓN GENERAL ───────────────────────────────────────────────────────

\section{Visión general}

Tu carta muestra esta distribución de elementos: Fuego 6.5, Tierra 1, Aire 3, Agua 6. El elemento más presente es Fuego (iniciativa, impulso y necesidad de movimiento) y Agua (sensibilidad, percepción emocional y conexión con el entorno). Esto señala un registro que suele aparecer con facilidad en tu forma de funcionar. El elemento menos presente es Tierra (estabilidad, concreción y contacto con lo práctico). No significa ausencia, sino una cualidad que puede necesitar más atención consciente. Predomina la modalidad mutable, asociada a la adaptación, la flexibilidad y el movimiento entre distintas situaciones. Tu Medio Cielo en Escorpio orienta tu presencia pública hacia la profundidad, intensidad y manejo de situaciones complejas.

\vspace{0.8cm}
\Needspace{5\baselineskip}


% ── LECTURA DE LOS ELEMENTOS ────────────────────────────────────────────

\section{Lectura de los elementos}

\subsection{Aire equilibrado}

Puedes pensar cuando lo necesitas y parar cuando no hace falta. La mente está disponible, pero no te arrastra.

\textbf{Observa:}
\begin{itemize}
\item cuándo usas la cabeza para resolver algo
\item cuándo sigues pensando sin necesidad
\item si puedes parar y descansar la mente
\end{itemize}

\vspace{0.4cm}
\Needspace{5\baselineskip}

\subsection{Tierra bajo}

Te cuesta mantener rutinas o seguir algo en el tiempo. Puedes empezar cosas pero no siempre sostenerlas.

\textbf{Observa:}
\begin{itemize}
\item si te cuesta mantener horarios o hábitos
\item si dejas cosas a medias
\item qué te ayudaría a tener más orden en tu día
\end{itemize}

\vspace{0.4cm}
\Needspace{5\baselineskip}

\subsection{Agua alto}

Sientes mucho. Lo que pasa por dentro tiene mucho peso. Puedes saturarte o quedarte enganchada en lo emocional.

\textbf{Observa:}
\begin{itemize}
\item cuándo te afecta demasiado lo que pasa
\item si te cuesta separar lo tuyo de lo de los demás
\item cómo te calmas cuando algo te desborda
\end{itemize}

\vspace{0.4cm}
\Needspace{5\baselineskip}

\subsection{Fuego alto}

Tienes mucha energía para actuar. Te mueves rápido y tomas iniciativa. Puedes ir demasiado deprisa o no sostener lo que empiezas.

\textbf{Observa:}
\begin{itemize}
\item si te lanzas sin pensar
\item si empiezas cosas y luego las dejas
\item cuándo necesitas bajar el ritmo
\end{itemize}

\vspace{0.4cm}
\Needspace{5\baselineskip}

\subsection{Integración}

Esto no define quién eres. Solo muestra cómo tiendes a funcionar.

No hay nada que cambiar. Se trata de verlo y aprender a manejarlo mejor.

\vspace{0.4cm}

\begin{center}
{\small\itshape
Ten en cuenta que tienes que hacer la amalgama de todos los elementos. Por ejemplo, si tienes Fuego muy alto pero también mucha Tierra, puede que no dejes proyectos a medias.
}
\end{center}


\vspace{0.8cm}
\Needspace{10\baselineskip}


% ── LOS TRES PILARES ────────────────────────────────────────────────────

\section{Los tres pilares}

\Needspace{14\baselineskip}
\subsection{Sol en Aries}

Con el Sol en Aries, necesitas movimiento, iniciativa y sensación de avance. Tu energía suele activarse rápidamente cuando algo te entusiasma o supone un reto. Muchas veces funcionas mejor actuando y corrigiendo sobre la marcha que esperando demasiado tiempo. Con el Sol en Casa 2, necesitas construir estabilidad y confianza a través de lo que haces, sostienes o desarrollas por ti. Los recursos, los valores personales y la sensación de autonomía suelen tener bastante peso en tu vida. Muchas veces necesitas sentir que puedes apoyarte en algo sólido construido desde ti.

\vspace{0.8cm}
\Needspace{5\baselineskip}

\Needspace{10\baselineskip}
\subsection{Luna en Sagitario}

Con la Luna en Sagitario, necesitas espacio, movimiento y sensación de amplitud para sentirte bien emocionalmente. Sueles recuperar equilibrio cuando puedes cambiar de perspectiva, aprender algo nuevo o abrir horizontes. La sensación de crecimiento y dirección suele influir bastante en tu estado emocional. Con la Luna en Casa 10, la vocación, los objetivos y el reconocimiento suelen influir bastante en tu estado emocional. Muchas veces lo profesional o lo visible hacia fuera tiene un impacto importante en cómo te sientes.

\vspace{0.8cm}
\Needspace{5\baselineskip}

\Needspace{10\baselineskip}
\subsection{Ascendente Acuario}

Con Ascendente en Acuario, normalmente te muestras de forma independiente, observadora y algo difícil de encajar en lo esperado. Sueles necesitar espacio y libertad para sentirte con comodidad en los entornos nuevos. Las demás personas pueden percibir cierta distancia o una forma distinta de estar y relacionarte.

\vspace{0.8cm}
\Needspace{5\baselineskip}

% ── CIERRE ───────────────────────────────────────────────────────────────

\section{Cierre}

Esta carta no busca definirte ni darte una identidad fija.

La astrología se utiliza aquí como lenguaje de observación:
una forma de mirar tendencias, ritmos y maneras de relacionarte con la vida.

La lectura ampliada desarrolla con más profundidad:
los vínculos,
la regulación emocional,
los patrones repetitivos,
la dirección de crecimiento,
las tensiones internas
y las configuraciones principales de la carta.

La lectura completa no busca darte más información.
Busca ayudarte a entender cómo sostener tu vida sin romperte por dentro.

\vspace{1cm}

\begin{center}

{\small\itshape\color{grisai}
Arquitectura Interna
}

\end{center}

\end{document}
